\documentclass[12pt]{article}
\usepackage[T1]{fontenc}
\usepackage[utf8]{inputenc}
\usepackage[brazil]{babel}
\usepackage[margin=1in]{geometry}
\usepackage{ragged2e}
\usepackage[normalem]{ulem}
\setlength{\parindent}{0pt}
\setlength{\parskip}{0.8\baselineskip}
\begin{document}
\justifying
\noindent Verifica-se que a questao objeto da controversia juridica suscitada pelo recurso manejado pela parte recorrente encontra-se afetada ao Tema de n. 344 da jurisprudencia da TNU - PU no processo n. PEDILEF 1006649-81.2020.4.01.3820/MG, afetado em 2023-10-19, voltado a elucidar a seguinte questao: "Salario-maternidade e devido em adocao de menor acima de doze anos?"

\noindent Por ocasião do julgamento, fixou-se a seguinte tese:

\noindent \textit{"E devido salario-maternidade por 120 dias ao adotante de menor de dezoito anos." (Tema 344 TNU).}

\noindent Contudo, verifica-se que, apesar de ter havido o julgamento do Tema de n. 344, em 2025-03-12, o respectivo Acórdão ainda não transitou em julgado.

\noindent Nesses termos, tendo em vista a irreversibilidade da medida nesta fase recursal, reputo mais adequado, em homenagem ao primado da segurança jurídica, aguardar o trânsito em julgado do referido Acórdão antes de autorizar a continuidade do processamento do recurso interposto nestes autos, que versa sobre a mesma controvérsia.

\noindent Nos termos do art. 16, §6º, VI, a e b,  do RITNU, as teses firmadas em representativo da controvérsia somente devem ser aplicadas após o trânsito em julgado do processo afetado. Confira-se:

\noindent \textit{Art. 16. Quando houver multiplicidade de recursos com fundamento em idêntica questão de direito, a Turma Nacional de Uniformização poderá afetar dois ou mais pedidos de uniformização de interpretação de lei federal como recurso representativo de controvérsia. [...] }\textit{\textbf{§ 6º O pedido de uniformização de interpretação de lei federal admitido como representativo da controvérsia será processado e julgado com observância deste procedimento: [...] VI - transitado em julgado o acórdão da Turma Nacional de Uniformizaçã}}\textit{o, os pedidos de uniformização de interpretação de lei federal sobrestados: a) terão seguimento denegado na hipótese de o acórdão recorrido coincidir com a orientação da Turma Nacional de Uniformização; ou b) serão encaminhados à Turma de origem para juízo de retratação, quando o acórdão recorrido divergir do decidido pela Turma Nacional, ficando integralmente prejudicados os pedidos de uniformização de interpretação de lei federal anteriormente interpostos.(Grifos acrescentados)}

\noindent Nesse contexto, determino o \textbf{SOBRESTAMENTO do PU} interposto nestes autos – o qual versa sobre a mesma matéria –, na forma do \uline{\textbf{art. 14, II, b)}},\textbf{ }c/c\textbf{ }\uline{\textbf{art. 16, § 6º, VI}}, ambos da Resolução CJF n. 586/2019, até o julgamento final do representativo da controvérsia, o PU no processo n. PEDILEF 1006649-81.2020.4.01.3820/MG (\textbf{Tema n. 344 da TNU}).

\noindent \textbf{Intimem-se. Cumpra-se.}
\end{document}
